% GENERATED FILE. Edit canonical lessons, then run npm run generate:books.
% canonical-source-hash: fnv1a64:73b199c6906bdf46
% canonical-lessons: ZH-C373-quedian, ZH-C373-haochu, ZH-C373-huaichu, ZH-C373-tianzhen, ZH-C373-dangzhen, ZH-R373-review-369-373, ZH-R373-review-final-373

\chapter{Pros and Cons}
\label{ch:zh-pros-and-cons}
\hlchaptermodality{373}

\begin{chapteropening}
\textbf{By the end of this chapter:} \emph{I can weigh up strong and weak points, benefits and harm, and tell someone not to take it seriously.}

Complete the last lesson of chapter 373: I can weigh up strong and weak points, benefits and harm, and tell someone not to take it seriously.
\end{chapteropening}

\section[quēdiǎn]{quēdiǎn — weak point; shortcoming}

\label{lesson:ZH-C373-quedian}



\begin{quote}
Before the new one: say the Mandarin for the purpose behind it, then the Mandarin for a strong point.
\end{quote}

\subsection*{You'll want to know: quēdiǎn}
\emph{quēdiǎn} — \textquotedblleft{}weak point; shortcoming\textquotedblright{}.

\begin{quote}
\emph{wǒ juéde měi ge rén dōu yǒu quēdiǎn} — I think everyone has their faults.
\end{quote}

Its characters come later; for now it is a word for the ear.

\begin{cousinweb}
\emph{quēdiǎn} is built from \emph{quē} and \zh{点} 'a point', both inherited from Middle Chinese, \emph{khjwiet} and \emph{temX} (Guangyun, in Baxter's transcription), which glosses the first 'too little' and, quoting an older dictionary, 'a broken vessel', and the second 'a dot'.
\end{cousinweb}

\begin{grammarlens}[title={What you've built}]
One more word for giving reasons.
\end{grammarlens}

\subsection*{Your turn}
\begin{itemize}
\raggedright
  \item \emph{Say it:} \emph{quēdiǎn}
  \item \emph{Say it:} \emph{quēdiǎn}, once more
  \item \emph{Say it:} \emph{wǒ juéde měi ge rén dōu yǒu quēdiǎn} — I think everyone has their faults.
  \item \emph{Recall:} say the Mandarin for the purpose behind it, then the Mandarin for a strong point, then say \emph{quēdiǎn} again
\end{itemize}

\begin{tcolorbox}[breakable,colback=teal!4,colframe=teal!35!black,title={Before you move on}]
What does quēdiǎn mean? (\textquotedblleft{}Weak point; shortcoming\textquotedblright{}.) Say it once more.
\end{tcolorbox}
\section[hǎochu]{hǎochu — benefit; advantage}

\label{lesson:ZH-C373-haochu}



\begin{quote}
Before the new one: say the Mandarin for a strong point, then the Mandarin for a weak point.
\end{quote}

\subsection*{You'll want to know: hǎochu}
\emph{hǎochu} — \textquotedblleft{}benefit; advantage\textquotedblright{}.

\begin{quote}
\emph{wǒ juéde měitiān yùndòng yǒu hǎochu} — I think exercising every day does you good.
\end{quote}

Its characters come later; for now it is a word for the ear.

\begin{cousinweb}
\emph{hǎochu} is built from \zh{好} 'good' and \emph{chù}, both inherited from Middle Chinese, \emph{xawX} and \emph{tsyhoH} (Guangyun, in Baxter's transcription), which glosses the first 'good' and the second 'a place'. The second is said light here.
\end{cousinweb}

\begin{grammarlens}[title={What you've built}]
One more word for giving reasons.
\end{grammarlens}

\subsection*{Your turn}
\begin{itemize}
\raggedright
  \item \emph{Say it:} \emph{hǎochu}
  \item \emph{Say it:} \emph{hǎochu}, once more
  \item \emph{Say it:} \emph{wǒ juéde měitiān yùndòng yǒu hǎochu} — I think exercising every day does you good.
  \item \emph{Recall:} say the Mandarin for a strong point, then the Mandarin for a weak point, then say \emph{hǎochu} again
\end{itemize}

\begin{tcolorbox}[breakable,colback=teal!4,colframe=teal!35!black,title={Before you move on}]
What does hǎochu mean? (\textquotedblleft{}Benefit; advantage\textquotedblright{}.) Say it once more.
\end{tcolorbox}
\section[huàichu]{huàichu — harm; downside}

\label{lesson:ZH-C373-huaichu}



\begin{quote}
Before the new one: say the Mandarin for a weak point, then the Mandarin for a benefit.
\end{quote}

\subsection*{You'll want to know: huàichu}
\emph{huàichu} — \textquotedblleft{}harm; downside\textquotedblright{}.

\begin{quote}
\emph{xīyān yǒu shénme huàichu} — What harm does smoking do?
\end{quote}

Its characters come later; for now it is a word for the ear.

\begin{cousinweb}
\emph{huàichu} is built from \emph{huài} and \emph{chù}, both inherited from Middle Chinese, \emph{hweajH} and \emph{tsyhoH} (Guangyun, in Baxter's transcription), which glosses the first 'to fall apart of itself' and the second 'a place'. The second is said light here.
\end{cousinweb}

\begin{grammarlens}[title={What you've built}]
One more word for giving reasons.
\end{grammarlens}

\subsection*{Your turn}
\begin{itemize}
\raggedright
  \item \emph{Say it:} \emph{huàichu}
  \item \emph{Say it:} \emph{huàichu}, once more
  \item \emph{Say it:} \emph{xīyān yǒu shénme huàichu} — What harm does smoking do?
  \item \emph{Recall:} say the Mandarin for a weak point, then the Mandarin for a benefit, then say \emph{huàichu} again
\end{itemize}

\begin{tcolorbox}[breakable,colback=teal!4,colframe=teal!35!black,title={Before you move on}]
What does huàichu mean? (\textquotedblleft{}Harm; downside\textquotedblright{}.) Say it once more.
\end{tcolorbox}
\section[\texorpdfstring{\zh{天真} (tiānzhēn)}{天真 (tiānzhēn)}]{\zh{天真} (tiānzhēn) — naive; innocent}

\label{lesson:ZH-C373-tianzhen}



\begin{quote}
Before the new one: say the Mandarin for a benefit, then the Mandarin for a downside.
\end{quote}

\subsection*{You'll want to know: \zh{天真}}
\textbf{\zh{天真}} — \emph{tiānzhēn} — \textquotedblleft{}naive; innocent\textquotedblright{}.

\begin{quote}
\textbf{\zh{虽然你说得有理但是你太天真了}} \emph{suīrán nǐ shuō de yǒulǐ dànshì nǐ tài tiānzhēn le} — You may have a point, but you're too naive.
\end{quote}

\begin{cousinweb}
\textbf{\zh{天真}} (\emph{tiānzhēn}) is built from \zh{天} 'the sky' and \zh{真} 'true', both inherited from Middle Chinese, \emph{then} and \emph{tsyin} (Guangyun, in Baxter's transcription), which glosses the first 'the dark heaven above' and the second 'true, as against false'. How the pair came to mean naive is not settled.
\end{cousinweb}

\begin{grammarlens}[title={What you've built}]
One more word for giving reasons.
\end{grammarlens}

\subsection*{Your turn}
\begin{itemize}
\raggedright
  \item \emph{Say it:} \emph{tiānzhēn}
  \item \emph{Say it:} \emph{tiānzhēn}, once more
  \item \emph{Say it:} \emph{suīrán nǐ shuō de yǒulǐ dànshì nǐ tài tiānzhēn le} — You may have a point, but you're too naive.
  \item \emph{Recall:} say the Mandarin for a benefit, then the Mandarin for a downside, then say \emph{tiānzhēn} again
\end{itemize}

\begin{tcolorbox}[breakable,colback=teal!4,colframe=teal!35!black,title={Before you move on}]
What does \zh{天真} mean? (\textquotedblleft{}Naive; innocent\textquotedblright{}.) Say it once more.
\end{tcolorbox}
\section[\texorpdfstring{\zh{当真} (dàngzhēn)}{当真 (dàngzhēn)}]{\zh{当真} (dàngzhēn) — to take it seriously; to take at face value}

\label{lesson:ZH-C373-dangzhen}



\begin{quote}
Before the new one: say the Mandarin for a downside, then the Mandarin for naive.
\end{quote}

\subsection*{You'll want to know: \zh{当真}}
\textbf{\zh{当真}} — \emph{dàngzhēn} — \textquotedblleft{}to take it seriously; to take at face value\textquotedblright{}.

\begin{quote}
\emph{nà shì qìhuà nǐ bié dàngzhēn} — Those were angry words; don't take them seriously.
\end{quote}

\begin{cousinweb}
\textbf{\zh{当真}} (\emph{dàngzhēn}) is built from \zh{当} and \zh{真} 'true', both inherited from Middle Chinese, \emph{tangH} and \emph{tsyin} (Guangyun, in Baxter's transcription), which glosses the first 'to be in charge' and the second 'true, as against false'.
\end{cousinweb}

\begin{grammarlens}[title={What you've built}]
One more word for giving reasons.
\end{grammarlens}

\subsection*{Your turn}
\begin{itemize}
\raggedright
  \item \emph{Say it:} \emph{dàngzhēn}
  \item \emph{Say it:} \emph{dàngzhēn}, once more
  \item \emph{Say it:} \emph{nà shì qìhuà nǐ bié dàngzhēn} — Those were angry words; don't take them seriously.
  \item \emph{Recall:} say the Mandarin for a downside, then the Mandarin for naive, then say \emph{dàngzhēn} again
\end{itemize}

\begin{tcolorbox}[breakable,colback=teal!4,colframe=teal!35!black,title={Before you move on}]
What does \zh{当真} mean? (\textquotedblleft{}To take it seriously; to take at face value\textquotedblright{}.) Say it once more.
\end{tcolorbox}
\section[(dialogue)]{Looking back: from Making It Clear to Pros and Cons}

\label{lesson:ZH-R373-review-369-373}



\begin{quote}
Say the Mandarin for to make something clear and however.
\end{quote}

\subsection*{Your turn}
\begin{itemize}
\raggedright
  \item \emph{Say it:} to make something clear, however, to give an example, secondly, in addition
  \item \emph{Say it:} a factor, all things considered, the crucial point, the effect something has, in short
  \item \emph{Say it:} even (going so far as), although, a conclusion, to sum up, a principle you live by
  \item \emph{Say it:} anyway, on earth (in a question), simultaneously, the purpose behind it, a strong point
  \item \emph{Say it:} a weak point, a benefit, a downside, naive, to take it seriously
  \item \emph{Read it:} \zh{但} \zh{虽} — each one aloud, then a word that holds it
\end{itemize}

\begin{tcolorbox}[breakable,colback=teal!4,colframe=teal!35!black,title={Before you move on}]
To make something clear? (\textbf{\zh{说明}}, \emph{shuōmíng}.) However? (\textbf{\zh{但是}}, \emph{dànshì}.) To give an example? (\emph{jǔlì}.) Secondly? (\emph{qícì}.) In addition? (\emph{lìngwài}.) A factor? (\emph{yīnsù}.) All things considered? (\emph{bìjìng}.) The crucial point? (\emph{guānjiàn}.) The effect something has? (\textbf{\zh{作用}}, \emph{zuòyòng}.) In short? (\emph{zǒngzhī}.) Even (going so far as)? (\emph{shènzhì}.) Although? (\textbf{\zh{虽然}}, \emph{suīrán}.) A conclusion? (\emph{jiélùn}.) To sum up? (\emph{zǒngjié}.) A principle you live by? (\emph{yuánzé}.) Anyway? (\emph{fǎnzhèng}.) On earth (in a question)? (\emph{dàodǐ}.) Simultaneously? (\textbf{\zh{同时}}, \emph{tóngshí}.) The purpose behind it? (\textbf{\zh{用意}}, \emph{yòngyì}.) A strong point? (\emph{yōudiǎn}.) A weak point? (\emph{quēdiǎn}.) A benefit? (\emph{hǎochu}.) A downside? (\emph{huàichu}.) Naive? (\textbf{\zh{天真}}, \emph{tiānzhēn}.) To take it seriously? (\textbf{\zh{当真}}, \emph{dàngzhēn}.)
\end{tcolorbox}
\section[(dialogue)]{Pros and Cons — once more}

\label{lesson:ZH-R373-review-final-373}



\begin{quote}
Say the Mandarin for a weak point and a benefit.
\end{quote}

\subsection*{Your turn}
\begin{itemize}
\raggedright
  \item \emph{Say it:} a weak point, a benefit, a downside, naive, to take it seriously
\end{itemize}

\begin{tcolorbox}[breakable,colback=teal!4,colframe=teal!35!black,title={Before you move on}]
A weak point? (\emph{quēdiǎn}.) A benefit? (\emph{hǎochu}.) A downside? (\emph{huàichu}.) Naive? (\textbf{\zh{天真}}, \emph{tiānzhēn}.) To take it seriously? (\textbf{\zh{当真}}, \emph{dàngzhēn}.)
\end{tcolorbox}
